\begin{table}[!htbp]
\centering
\footnotesize
\caption{Continuous time-respecting film-visibility model.}
\label{tab:si-visibility-models}
\renewcommand{\arraystretch}{1.05}
\setlength{\tabcolsep}{3pt}
\begin{tabular}{>{\raggedright\arraybackslash}p{0.22\linewidth}>{\raggedright\arraybackslash}p{0.20\linewidth}>{\centering\arraybackslash}p{0.082\linewidth}>{\centering\arraybackslash}p{0.07\linewidth}>{\centering\arraybackslash}p{0.095\linewidth}>{\centering\arraybackslash}p{0.06\linewidth}>{\centering\arraybackslash}p{0.06\linewidth}}
\toprule
Model & Term & Estimate & \makecell[b]{Cluster\\ SE} & 95\% CI & N obs. & N songs \\
\midrule
Film visibility (continuous, time-respecting, primary) & log1p max recorded worldwide box office; films released by snapshot year & 1.15 & 0.102 & [0.95, 1.35] & 15,245 & 4,294 \\
\bottomrule
\end{tabular}
\begin{flushleft}\footnotesize
Song-snapshots linked to films released by the snapshot year with recorded worldwide box office. Recorded totals may include receipts after the snapshot date. The model adjusts for age, squared age, Billboard weeks, peak position, realized film count, and snapshot fixed effects, with cluster-by-song intervals. Box office measures film-level visibility, not scene-level soundtrack prominence.
\end{flushleft}
\end{table}
